\documentclass[
  bibliography=totoc,     % Literatur im Inhaltsverzeichnis
  captions=tableheading,  % Tabellenüberschriften
  titlepage=firstiscover, % Titelseite ist Deckblatt
  parskip=half,
  ngerman
]{scrartcl}

% Paket float verbessern
\usepackage{scrhack}

% Warnung, falls nochmal kompiliert werden muss
\usepackage[aux]{rerunfilecheck}

% unverzichtbare Mathe-Befehle
\usepackage{amsmath}
% viele Mathe-Symbole
\usepackage{amssymb}
% Erweiterungen für amsmath
\usepackage{mathtools}

% Fonteinstellungen

% Latin Modern Fonts werden automatisch geladen
% Alternativ zum Beispiel:
%\setromanfont{Libertinus Serif}
%\setsansfont{Libertinus Sans}
%\setmonofont{Libertinus Mono}

% Wenn man andere Schriftarten gesetzt hat,
% sollte man das Seiten-Layout neu berechnen lassen
\recalctypearea{}

% deutsche Spracheinstellungen
\usepackage[main=ngerman]{babel}

%----------------
% Mathematik-Stil nach ISO
%----------------
\usepackage[
  math-style=ISO,
  bold-style=ISO,
  sans-style=italic,
  nabla=upright,
  partial=upright,
  warnings-off={
    mathtools-colon,
    mathtools-overbracket,
  },
]{unicode-math}

% traditionelle Fonts für Mathematik
\setmathfont{Latin Modern Math}
\setmathfont{XITS Math}[range={scr, bfscr}]
\setmathfont{XITS Math}[range={cal, bfcal}, StylisticSet=1]

% Zahlen und Einheiten
\usepackage[
  locale=DE,
  separate-uncertainty=true,
  per-mode=symbol-or-fraction,
]{siunitx}

% chemische Formeln
\usepackage[
  version=4,
  math-greek=default,
  text-greek=default,
]{mhchem}

% richtige Anführungszeichen
\usepackage[autostyle]{csquotes}

% schöne Brüche im Text
\usepackage{xfrac}

% Standardplatzierung für Floats einstellen
\usepackage{float}
\floatplacement{figure}{htbp}
\floatplacement{table}{htbp}

% Floats innerhalb einer Section halten
\usepackage[
  section,
]{placeins}
\makeatletter
\AtBeginDocument{%
  \expandafter\renewcommand\expandafter\subsection\expandafter{%
    \expandafter\@fb@secFB\subsection
  }%
}
\makeatother

% Seite drehen für breite Tabellen: landscape Umgebung
\usepackage{pdflscape}

% Seitenränder temporär kleiner machen
\usepackage{geometry}

% Captions schöner machen.
\usepackage[
  labelfont=bf,
  font=small,
  width=0.9\textwidth,
]{caption}
% subfigure, subtable, subref
\usepackage{subcaption}

% Grafiken können eingebunden werden
\usepackage{graphicx}
\usepackage{grffile}
\graphicspath{{./bilder/}{./}}


% schöne Tabellen
\usepackage{booktabs}

% Verbesserungen am Schriftbild
\usepackage{microtype}

% Literaturverzeichnis
\usepackage[
  backend=biber,
  natbib=true,
  style=numeric,
  sorting=none
]{biblatex}
\addbibresource{../bib/lit.bib}
\addbibresource{../bib/programme.bib}
\addbibresource{../bib/versuche.bib}

% Hyperlinks im Dokument
\usepackage[
  german,
  unicode,
  pdfusetitle,
  pdfcreator={},
  pdfproducer={},
]{hyperref}
\usepackage{hyperref}

% erweiterte Bookmarks im PDF
\usepackage{bookmark}

% Trennung von Wörtern mit Strichen
\usepackage[shortcuts]{extdash}

% Anpassbare Enumerates/Itemizes
\usepackage{enumitem}

% Tabellen mit einer Zelle über mehrere Reihen
\usepackage{multirow}

% Tabellen aus .csv Dateien
\usepackage{csvsimple}

%----------------
% Neue Commands
%----------------
\NewDocumentCommand \dif{m}{
  \mathinner{\symup{d} #1}
}
\NewDocumentCommand \tableSI{m m}{
  $#1 \:/\: \si{#2}$
}

%----------------
% Autoreninformationen
%----------------

\author{%
  Jannes Krause\\
  \href{mailto:jannes.krause@tu-dortmund.de}{jannes.krause@tu-dortmund.de}%
  \and
  Lea Krüger\\
  \href{mailto:lea.krueger@tu-dortmund.de}{lea.krueger@tu-dortmund.de}%
}





